An \emph{account document} describes the owner's policy, including
administrative rules and recovery; a \emph{grant document} describes a bounded
set of rules for one application. Applying a grant document creates or replaces
an \emph{installed grant}, the rules the account consults, stored under a grant
ID and stamped with the current grant generation. Grant~1 as the reference run
installs it, abbreviated:
\def\vpass{\fact{standalone-passphrase}}\def\vlimit{\fact{grant1-cap-stroops}}%
\def\vperiod{\fact{grant1-period-ledgers}}%
\begin{Verbatim}[fontsize=\scriptsize,xleftmargin=1em,commandchars=\\\(\)]
{ "version": 1, "network": "\vpass",
  "signers": [ { "id": "app", "verifier": "C...ED25519", "key": "...Ed25519 key, hex..." } ],
  "rules": [ { "name": "spend", "scope": { "type": "contract", "address": "C...XLM" },
               "principals": { "type": "all", "signers": ["app"] }, "functions": ["transfer"],
               "cap": { "limit": "\vlimit", "period-ledgers": \vperiod } } ] }
\end{Verbatim}
\code{scope} is the rule target and \code{principals} the signer condition;
\code{cap.limit} is \fact{grant1-cap-xlm}~XLM in stroops, charged against the
amount argument of \code{transfer(from, to, amount)}, the only function a capped
rule may list. A rule names one target contract, a function list, argument
conditions and a threshold over listed signers, and optionally an expiry and a
cap. Its \fact{doc-conditions-count} kinds of argument condition,
\fact{doc-conditions}, are conjunctive, and a document has no negation. A
document cannot bound a single call's amount: a capped rule gets only a type
guard that argument~2 is an \code{i128}.

\paragraph{Installation.} The owner authorizes \code{apply\_grant} (exported as
\code{apply\_scope}; Table~\ref{tab:tab-api}). The stateless document compiler
validates the JSON, resolves verifier names through the \code{perch/verifiers}
registry and returns canonical bytes, their hash \code{doc\_hash} and a model.
The account lowers the model with rule-core's \code{lower\_scoped} under its
address, its excluded set and the current generation, the step the lowering and
containment theorems cover. Grant lowering refuses (\code{ScopeEscalation}) a
document that configures recovery, and any rule whose target is the account, its
adopted recovery controller, or the zero-knowledge adapter and pool of its
recovery configuration. The account also refuses non-normal documents, rules with
no signer, and credentials that recovery revoked or a stale grant named, then
stores the grant; a client can rebuild its JSON from chain state and check it
against \code{doc\_hash} (Appendix~\ref{sec:formal-appendix}).

\paragraph{Rule semantics.} A rule set denotes the union of its rules'
permissions, but the account does not search it: the transaction names one rule
per context, which must itself allow the context, be unexpired and, if it
carries a generation, be current. A supplied signer in no chosen rule is
refused. rule-core's \fact{attenuation-steps} attenuation steps each return a rule
set allowing nothing the original does not (\ref{stmt:attenuation-narrows-authority},
\ref{stmt:trace-inclusion}; Appendix~\ref{sec:formal-appendix}); no account entry
point installs an attenuated grant at the pin.

\paragraph{Sliding windows.} For period $P$, the window at ledger $L$ holds the
charges made after $L-P$ and through $L$ (Fig.~\ref{fig:fig-cap-window}). The
evaluator checks the new total against the window and returns a new record;
contexts selecting one rule in one authorization see each other's charges, and a
negative or malformed amount, overflow, a ledger before the newest charge and a
full history fail closed. The \fact{ledger-entry-limit-bytes}-byte entry limit
holds \fact{spends-per-window} spends per record, so a grantee can exhaust its
own window with small transfers.
